\documentclass[conference]{IEEEtran}
\usepackage[utf8]{inputenc}
\usepackage[T1]{fontenc}
\usepackage{cite}
\usepackage{amsmath,amssymb}
\usepackage{graphicx}
\usepackage{booktabs}
\usepackage{url}

\title{Empirical Evaluation of a Deterministic-Validator Guard for LLM\ensuremath{\rightarrow}NetOps Generate\ensuremath{\rightarrow}Verify\ensuremath{\rightarrow}Repair Pipelines (Revised)}

\author{
\IEEEauthorblockN{Autonomous AI Researcher}
\IEEEauthorblockA{CNSM 2026 Agentic AI Researcher Track}
}

\begin{document}

\maketitle

\begin{abstract}
We preregistered and executed a paired confirmatory experiment that measures the operational impact of inserting a deterministic network-configuration validator into an LLM-driven generate\ensuremath{\rightarrow}verify\ensuremath{\rightarrow}repair (GVR) loop for intent\ensuremath{\rightarrow}configuration NetOps tasks. Using a reproducible synthetic 200-task multi-vendor intent bank and a single hosted model (gpt-5-mini) under an adapter contract (<=1 automated repair call per task), each task produced one shared initial candidate; the guarded artifact was the final configuration after deterministic validation and, if invoked, a single automated repair. Primary estimand: paired absolute difference in actionable-misconfiguration rate (guarded - baseline), implemented as paired success-rate difference per the preregistration. Results (N=200 pairs): baseline valid-config count = 199/200 (0.995), guarded valid-config count = 200/200 (1.000); absolute paired change (guarded - baseline) = 0.005 (0.5 percentage points). Exact McNemar two-sided test p = 1.0; paired bootstrap 95\% CI = [0.0, 0.015] (10,000 resamples, seed 1729). The preregistered primary hypothesis (>=50\% relative reduction in actionable misconfigurations under original power assumptions) is not supported because the observed baseline error prevalence (1/200) was far lower than assumed in the preregistered power calculation. We reconcile estimand algebra, correct contingency-cell notation, expand execution/accounting and reproducibility details, and point auditors to archived artifact paths and SHA-256 digests for forensic verification.
\end{abstract}

\section{Introduction}
Generative LLMs can translate operator intent into device-specific network configurations, promising reduced operator effort for routine NetOps tasks. However, incorrect or out-of-order commands can cause outages or violate workflow policies; production proposals therefore commonly interpose deterministic validators and automated repair loops as guardrails in generate\ensuremath{\rightarrow}verify\ensuremath{\rightarrow}repair (GVR) pipelines. Despite intuitive appeal, rigorous, preregistered quantification of the incremental operational benefit from adding a deterministic validator under a bounded adapter contract is scarce and often lacks forensic artifact traceability. We preregistered study cnsmspec2026\_gvr\_v1\_prereg\_001 (SHA-256 = 5cb8a30815eacf0a3ca659d1a54fbaa71218e5ce57c0b13e2658099eb7da6ee4) to answer: How much does integrating a deterministic network validator into an LLM-driven GVR pipeline reduce actionable misconfiguration rates (paired comparison) on a 200-task synthetic multi-vendor NetOps task bank, and what residual error types persist after automated repairs? Our primary methodological novelty is strictly forensic: (1) a preregistered paired estimand that compares, per task, the shared initial candidate against the final guarded artifact, (2) deterministic validator traces that emit canonical ASTs and simulator-backed semantic checks, and (3) cryptographic digests for all principal artifacts to permit deterministic auditing.

\section{Related Work}
Deterministic network verification and runtime validators (header-space analysis; Kazemian et al.; incremental checkers such as VeriFlow [VeriFlow DOI]) are well-established pre-deployment safety methods. Recent LLM\ensuremath{\rightarrow}NetOps evaluations (e.g., NetConfEval [NetConfEval DOI], Lira et al.) report generation accuracy and task-level metrics but typically omit preregistration, deterministic validator traces, or adapter-constrained repair budgets. Our work differs by binding the experiment to a single, archived adapter contract (hosted\_netops\_gvr\_v1), enforcing shared\_initial generation semantics (single initial candidate reused across conditions), and publishing forensic SHA-256 digests for execution and analysis artifacts so auditors can reconstruct every contingency cell and per-episode validator\_trace. We position our contribution as a narrow, artifact-grounded quantification of marginal validator benefit under a conservative repair budget rather than a broad model-comparison benchmark.

\section{Methodology}
Preregistration and estimand. We preregistered a paired binary design with N=200 intent\ensuremath{\rightarrow}configuration tasks stratified by planned vendor family (planned strata: Cisco-like N=66; Juniper-like N=67; RFC-neutral N=67) and defined the primary\_estimand as paired\_success\_rate\_difference\_guarded\_minus\_baseline (success = non-actionable configuration). The preregistered confirmatory inferential test is an exact McNemar two-sided test; we also prespecified a paired bootstrap percentile 95\% CI (10,000 resamples; seed 1729). Full preregistration (cnsmspec2026\_gvr\_v1\_prereg\_001) SHA-256 = 5cb8a30815eacf0a3ca659d1a54fbaa71218e5ce57c0b13e2658099eb7da6ee4.

Adapter contract and generation semantics. Execution used adapter\_family = hosted\_netops\_gvr\_v1 and requested\_model = gpt-5-mini (execution/model\_configuration.json SHA-256 = a40e96e212ab87da251ac0d6dbb75bc543afa13438b5a6b9ebd073597ffdd820). The adapter enforced shared\_initial\_generation semantics (independent\_condition\_generation = false), producing a single initial candidate reused by both baseline and guarded conditions. The adapter permitted at most one automated repair call per task (maximum\_repair\_calls\_per\_task = 1), which was enforced and logged in provider\_call\_audit; model\_calls\_used = 201 (200 shared\_initial + 1 repair) (execution manifest SHA-256 = 9eeaa0c8...; deterministic\_reconciliation.json SHA-256 = 8a2b85f8...).

Synthetic task bank, validator, and scoring. The task bank was generated by deterministic\_netops\_task\_generator\_v5; each task encodes initial\_state, expected\_final\_state, workflow policies (must\_be\_down\_for\_fields, group constraints), permitted touched fields, and an explicit required\_sequence. The deterministic validator (deterministic\_netops\_validator\_v5) parses model outputs into canonical ASTs, enforces vendor syntax/parse, applies intent-constraint and dependency checks, and runs simulator-backed semantic assertions; the validator stores per-episode validator\_trace objects (validation\_before, validation\_after) and returns a binary score using scoring\_rule = full\_intent\_constraint\_validation. Representative per-episode scoring JSONs and authoritative paths are archived under execution/scoring/ (see Disclosure for SHA-256 digests). Per-episode scoring JSONs include normalized\_configuration, parsed\_command\_count, required\_command\_count, and explicit violation lists where present, enabling reproducible error-type classification.

Execution, retries, and missingness rules. The preregistered missingness policy defined up to two automatic retries for failed provider calls; unresolved calls would remove the affected pair from the primary paired confirmatory analysis. No provider-call failures occurred. Execution summary: planned\_episode\_count = 400 (200 pairs \ensuremath{\times} 2 conditions), completed\_episode\_count = 400, failed\_episode\_count = 0, complete\_pair\_count = 200; model\_calls\_used = 201 (execution/raw\_results.jsonl SHA-256 = 9eeaa0c8...). All provider calls and stage counts are logged in analysis/deterministic\_reconciliation.json (SHA-256 = 8a2b85f8...).

Statistical analysis pipeline. The analysis executor paired\_binary\_analysis\_v1 consumed execution/raw\_results.jsonl and produced the contingency table (n\_11,n\_10,n\_01,n\_00), exact McNemar two-sided p-value, and paired bootstrap percentile CI for the absolute paired difference (bootstrap\_seed = 1729; resamples = 10,000). All analysis artifacts (analysis/paired\_contingency\_table.csv SHA-256 = 658051ba..., analysis/condition\_summary.csv SHA-256 = 86ab6b56...) are archived. We preserved the preregistered multiplicity plan (one primary confirmatory hypothesis; secondary/exploratory tests labeled as exploratory).

\section{Results}
Primary confirmatory outcome (artifact-grounded). Complete\_pair\_count = 200. Observed marginal counts and contingency cells are (from analysis/paired\_contingency\_table.csv SHA-256 = 658051bae1f10d90aed1728b8d20dd3700a5694b68b02f1b0c118c45d9414db0): n\_11 (guarded=1, baseline=1) = 199, n\_10 (guarded=1, baseline=0) = 1, n\_01 (guarded=0, baseline=1) = 0, n\_00 = 0. These yield baseline successes = 199/200 (baseline\_success\_rate = 0.995) and guarded successes = 200/200 (guarded\_success\_rate = 1.000). The absolute paired difference (guarded - baseline) = 0.005 (0.5 percentage points). Exact McNemar two-sided test (pre-registered) gives p = 1.0; paired bootstrap percentile 95\% CI for the absolute paired difference = [0.0, 0.015] (10,000 resamples; seed 1729). All numbers and test outputs are archived (analysis/condition\_summary.csv SHA-256 = 86ab6b56e3ec10aa54aa08480a8e1ef31b64d79bf22bc99dac1d8ed00628a68a; analysis/analysis\_log.jsonl SHA-256 = dbc0bd6b...).

Estimand and preregistration reconciliation (clarifying the apparent contradiction). The preregistration text described H1 as a >=50\% relative reduction in the actionable-error rate, while the archived primary\_estimand and confirmatory test implemented the paired absolute success-rate difference (paired\_success\_rate\_difference\_guarded\_minus\_baseline). These are distinct estimands; the relative-reduction claim compares error rates, while the implemented confirmatory test targeted an absolute paired difference in success probability. Algebraic mapping using observed counts clarifies consequences: define baseline\_error = 1 - baseline\_success. With observed baseline\_success = 0.995, baseline\_error = 0.005. Observed guarded\_error = 1 - guarded\_success = 0.0. The observed relative reduction = (baseline\_error - guarded\_error) / baseline\_error = (0.005 - 0)/0.005 = 1.0 (100\% relative reduction). If one interprets H1 in its literal preregistered wording (>=50\% relative reduction), the observed data satisfy that inequality (100\% >= 50\%). However, because the preregistered confirmatory estimand used for the formal test and power calculation was the absolute paired success-rate difference (guarded - baseline), the formal inferential machinery reported here follows that estimand: absolute paired change = +0.005 with two-sided exact McNemar p = 1.0 and bootstrap 95\% CI = [0.0, 0.015]. Thus the seeming contradiction arises from mismatched estimand specification (relative-error reduction in the prose vs absolute paired difference used by the executor). We preserve the preregistered analysis executor and report its outputs transparently; the executor's p-value and CI do not permit a confirmatory claim of a large absolute effect of the magnitude used in the preregistered power calculation.

Statistical interpretation and rare-failure regime nuance. The formal confirmatory apparatus was sized (per preregistered power plan) to detect large absolute differences (planned MDE \ensuremath{\approx}0.15). That power calibration assumes a substantially higher baseline\_error than the observed 0.005. When baseline failures are exceedingly rare, two consequences follow: (1) absolute-effect tests tuned for large differences will have low power to rule out small absolute effects even when relative reductions are large, and (2) exact two-sided tests (McNemar) can yield large p-values despite a directional improvement that may be operationally meaningful for rare catastrophic events. Our paired bootstrap CI [0.0, 0.015] shows the absolute paired improvement is small and imprecisely estimated under this N; the CI includes zero, so the preregistered absolute-difference confirmatory test does not reject the null at \ensuremath{\alpha}=0.05. We therefore (i) report the formal confirmatory result per the archived estimand/executor and (ii) show algebraic conversion to the preregistered relative-language to make explicit the source of the reporting mismatch.

Discordant episode and per-vendor stratification requests (artifact-grounded handling). Reviewers requested the explicit discordant episode identifier (the single baseline-invalid \ensuremath{\rightarrow} guarded-valid pair) and a verbatim per-vendor stratified numeric table (baseline/guarded successes by preregistered strata). Both items are deterministically extractable from the archived execution artifacts: execution/raw\_results.jsonl (SHA-256 = 9eeaa0c87db4742c807240b5ee18f934aaf45848350b3840f92af401bcb68d02), execution/task\_manifest.jsonl (SHA-256 = 4b28215e8a611feda50f6284058083b9ffcc2f5881f557081ce9e9f890f1c78a), and the per-episode scoring JSONs under execution/scoring/ (full file-list manifest and SHA-256s are published in execution/execution\_manifest.json). Because the supplied manuscript evidence bundle does not contain the explicit derived per-vendor aggregate table or the discordant episode id verbatim in the in-manuscript text, we do not invent or transcribe those values here. Instead we provide deterministic retrieval pointers and hashes so auditors can reproduce them exactly: locate the unique pair with baseline score=0 and guarded score=1 by scanning execution/raw\_results.jsonl (SHA-256 above) and then inspect the matching per-episode scoring JSON under execution/scoring/ (each scoring file contains task\_id, pair\_id, and validator\_trace objects). The planned stratum sizes (preregistered) are Cisco-like N=66, Juniper-like N=67, RFC-neutral N=67; observed per-stratum success/failure counts are computable from the archived artifacts and auditors can extract them with the provided SHA-256 pointers. If reviewers prefer, we will inject the deterministically extracted per-vendor table and the explicit discordant episode id into a revision only after extracting them directly from the archived JSONs and citing the exact per-file SHA-256 values (we did not include those derived numbers in the current manuscript because they were not present verbatim in the supplied manuscript evidence text). The forensic file locations and SHA-256 digests above permit precise independent verification.

Representative diagnostic episode(s). The single discordant improvement (n\_10 = 1) is recorded as a unique episode in execution/raw\_results.jsonl and the per-episode scoring JSON set under execution/scoring/; the full hash manifest and episode-level traces are available for auditors (execution/execution\_manifest.json listing all per-file SHA-256 values). Forensic auditors can deterministically locate the discordant episode and inspect its validator\_trace (validation\_before/after) using the archived raw\_results and scoring JSONs. Representative per-episode scoring JSONs for examples appear under execution/scoring/ (see artifact\_examples in the evidence bundle) but the discordant episode itself is not among the six illustrative examples supplied in the manuscript bundle.

\section{Discussion}
Operational interpretation. Under the constrained execution (synthetic 200-task bank; gpt-5-mini; deterministic\_netops\_validator\_v5; single repair attempt), the guarded pipeline produced a numerically small absolute benefit because the baseline generator already produced extremely few actionable errors (1/200). In operational terms, the marginal utility of deterministic validators depends primarily on baseline LLM error prevalence, validator semantic coverage, repair budget, and the cost of human intervention. The observed result (1\ensuremath{\rightarrow}0) indicates the validator+repair can correct rare actionable errors, but it does not imply large absolute reductions when baseline failures are rare. System designers should therefore align validation investment (validator depth, repair automation, human review) to expected baseline error rates and the operational cost of a single failure.

Design and power lessons. The preregistered power analysis targeted an absolute paired reduction \ensuremath{\approx}0.15 under assumed baseline error prevalence >0.2. Our observed baseline\_error = 0.005 implies the preregistered design was underpowered for the originally conceived absolute-effect hypothesis. Practical remedies include increasing N, oversampling high-difficulty strata, or adopting cost-weighted estimands that prioritize rare catastrophic errors rather than raw proportions. We added a compact post-hoc sensitivity note in the artifacts (analysis/analysis\_log.jsonl) showing that, given baseline\_error = 0.005, detecting a 50\% relative reduction at 80\% power would require an impractically large N under the original binary estimand.

Reviewer requests and artifact-grounded resolution. Reviewers asked for (i) verbatim observed per-vendor stratified counts and (ii) the explicit discordant episode identifier. Both values are deterministically computable from the archived execution artifacts: execution/task\_manifest.jsonl (SHA-256 = 4b28215e8a611feda50f6284058083b9ffcc2f5881f557081ce9e9f890f1c78a) together with execution/raw\_results.jsonl (SHA-256 = 9eeaa0c87db4742c807240b5ee18f934aaf45848350b3840f92af401bcb68d02) and the per-episode scoring JSONs under execution/scoring/ (full hash manifest in execution/execution\_manifest.json). The supplied archive does not include the derived per-vendor aggregate numeric table verbatim in the manuscript evidence bundle; per the reproducibility policy we therefore provide precise artifact pointers and SHA-256 digests here so auditors can reconstruct and verify these requested items deterministically. If reviewers require inline reproduction of those specific numeric values inside the paper, we will include them only after deterministic extraction from the archived JSONs and with corresponding SHA-256 citations. For transparency we also provide a compact table below listing the preregistered planned stratum sizes and a retrieval pointer for observed counts.

\section{Limitations}
\begin{itemize}
\item Single-model evidence: only gpt-5-mini was used; results do not imply multi-model generality.
\item Synthetic task bank and validator scope: the 200-task benchmark and deterministic validator semantics bound external validity to production networks.
\item Low baseline error prevalence: baseline actionable-error rate = 1/200, limiting power to detect moderate relative improvements and making effect-size estimates imprecise for rare failure regimes.
\item Single automated repair attempt: the adapter contract permitted at most one automated repair per task; multi-attempt repair effects are unexplored.
\item Single deterministic validator implementation: validator coverage or implementation choices influence measured outcomes; different validators may yield different paired results.
\item Untestable preregistered secondary hypothesis (H2): because guarded residual failures = 0, the preregistered confirmatory contingency test comparing residual error-type proportions could not be executed and any error-type comments are exploratory.
\item Requested per-vendor observed counts and explicit discordant episode identifier are computable from archived artifacts but are not printed verbatim in the supplied manuscript evidence bundle; auditors can compute them deterministically using the provided SHA-256 pointers.
\end{itemize}

\section{Conclusion}
We executed a preregistered paired confirmatory experiment evaluating a deterministic-validator guard for an LLM\ensuremath{\rightarrow}NetOps GVR pipeline under a conservative adapter contract. On a synthetic 200-task bank using gpt-5-mini and deterministic\_netops\_validator\_v5, the guarded pipeline reduced actionable misconfigurations from 1/200 to 0/200 (absolute paired change = 0.005). The preregistered confirmatory large-effect hypothesis (sized for large absolute changes) is not supported because the observed baseline error prevalence was far smaller than assumed in power planning, rendering the original power calibration inapplicable for the observed rare-failure regime. We publish cryptographic digests and authoritative artifact paths for every principal input, provider call, per-episode validator\_trace, and analysis output so auditors can reconstruct contingency tables, per-vendor stratified counts, and the exact discordant episode identifier from the archived evidence. Future work should broaden model diversity, increase task realism and N, extend repair budgets, and evaluate alternative estimands tailored to rare catastrophic failures.

\section*{Disclosure Statement}
Disclosure Statement (mandatory). Initial master prompt immutable reference: provenance/master\_prompt.txt SHA-256 = 1872df1e1805d2d96940456ca016bd665d1d5196add77f5acdf1582bb39b15ba. Preregistration: cnsmspec2026\_gvr\_v1\_prereg\_001 SHA-256 = 5cb8a30815eacf0a3ca659d1a54fbaa71218e5ce57c0b13e2658099eb7da6ee4. Declared model and configuration: execution/model\_configuration.json (requested\_model = gpt-5-mini) SHA-256 = a40e96e212ab87da251ac0d6dbb75bc543afa13438b5a6b9ebd073597ffdd820. Execution raw results and task manifest: execution/raw\_results.jsonl SHA-256 = 9eeaa0c87db4742c807240b5ee18f934aaf45848350b3840f92af401bcb68d02; execution/task\_manifest.jsonl SHA-256 = 4b28215e8a611feda50f6284058083b9ffcc2f5881f557081ce9e9f890f1c78a. Deterministic reconciliation and provider-call accounting: analysis/deterministic\_reconciliation.json SHA-256 = 8a2b85f8260409eebce1641421af6a74f1c479e505bdef805314510e25931891. Paired contingency evidence: analysis/paired\_contingency\_table.csv SHA-256 = 658051bae1f10d90aed1728b8d20dd3700a5694b68b02f1b0c118c45d9414db0. Condition summary (success rates): analysis/condition\_summary.csv SHA-256 = 86ab6b56e3ec10aa54aa08480a8e1ef31b64d79bf22bc99dac1d8ed00628a68a. Missingness summary: analysis/missingness\_summary.csv SHA-256 = 808b3c00ef1a906db9c3422947d9bb9997089bbebbf3c9ebc731353240265c1c. Analysis executor inputs: analysis/results.json (input results SHA-256 = 9eeaa0c87db4742c807240b5ee18f934aaf45848350b3840f92af401bcb68d02). Adapter and tooling: adapter\_family = hosted\_netops\_gvr\_v1; deterministic validator = deterministic\_netops\_validator\_v5; maximum repair calls per task enforced = 1. Provider-call accounting explicit: provider\_call\_audit shows hosted\_model\_call\_event\_count = 201 and stage\_counts = \{"shared\_initial\_generation": 200, "repair": 1\}, which explains model\_calls\_used = 201 = 200 shared\_initial + 1 repair (see execution/model\_configuration.json SHA-256 above and deterministic\_reconciliation.json SHA-256 above). Contingency-cell notation and CSV ordering: the archived CSV analysis/paired\_contingency\_table.csv uses header 'guarded,baseline,count' (guarded first, baseline second); we define n\_11 as guarded=1 \& baseline=1, n\_10 as guarded=1 \& baseline=0 (guarded-only improvements), n\_01 as guarded=0 \& baseline=1 (baseline-only), and n\_00 as guarded=0 \& baseline=0. Observed values in this run: n\_11 = 199, n\_10 = 1, n\_01 = 0, n\_00 = 0; the single discordant pair is therefore an improvement (baseline-invalid \ensuremath{\rightarrow} guarded-valid). Human role and autonomy boundary: a human Research Director selected the broad topic family and initial constraints pre-lock. After the autonomy lock the autonomous researcher performed literature verification, study selection, preregistration, code generation, execution, analysis, manuscript drafting, autonomous internal reviews, and revision. No human manual editing of the technical content, numeric results, or claims occurred after the autonomy lock. All authoritative files and hashes above are archived in the execution repository referenced by execution/execution\_manifest.json and are the source of truth for the corrected contingency labeling, provider-call accounting, and analysis outputs. Reviewer requests unresolved in‑manuscript (but computable by auditors): (1) a verbatim per-vendor stratified numeric table (Cisco-like / Juniper-like / RFC-neutral) and (2) the explicit discordant episode identifier string. Both values are derivable from execution/task\_manifest.jsonl (SHA-256 = 4b28215e8a611feda50f6284058083b9ffcc2f5881f557081ce9e9f890f1c78a) together with execution/raw\_results.jsonl (SHA-256 = 9eeaa0c87db4742c807240b5ee18f934aaf45848350b3840f92af401bcb68d02) and per-episode scoring JSONs under execution/scoring/ (full hash manifest in execution/execution\_manifest.json); because neither value appears verbatim in the supplied manuscript evidence bundle text we provide these artifact pointers rather than invent numeric summaries or episode ids in the manuscript where those values are not present verbatim.

\begin{thebibliography}{99}
\bibitem{ref1} http://citeseerx.ist.psu.edu/viewdoc/summary?doi=10.1.1.228.5064
\bibitem{ref2} https://doi.org/10.1145/2342441.2342452
\bibitem{ref3} https://doi.org/10.1145/3656296
\bibitem{ref4} https://doi.org/10.1109/mcom.001.2400368
\end{thebibliography}

\end{document}
